\section{Global growth and cone geometry}
\label{sec:conclusion}

All parameters, the finite degree range, and the exact sequence of controls
are now fixed. We first prove that the selected center-regular solutions
satisfy the pointwise bound in Theorem~\ref{thm:main}. We then examine the
cone limits, which explain why additional hypotheses in the comparison
literature do not apply.

\subsection{Every-point growth and the strict count}
\begin{proof}[Proof of Theorem~\ref{thm:main}]
We convert the exact cycles from Proposition~\ref{prop:transmission} into polynomial growth.
For each \(l\) with \(M_l>0\) and each \(1\leq i\leq M_l\),
choose the center-regular solution with value vector
\(Y(t_{j_0})=e_{l,i}\).  Lemma~\ref{lem:regular} permits this
choice.  Near the origin it is \(r^l\) times an element of
\(V_l\), hence a homogeneous harmonic polynomial in Euclidean
coordinates.  It is smooth and harmonic there and, by
\eqref{eq:harmonic-ode}, on the rest of the manifold.

Let \(\sigma_l\) be the cyclic permutation of
\(\{1,\ldots,M_l\}\) underlying \(\Pi_l\), and define
\(i_{j_0}=i\), \(i_{j+1}=\sigma_l(i_j)\).  Applying
\eqref{eq:exact-transmission} sends \(e_{l,i_j}\) to a signed
multiple of \(e_{l,i_{j+1}}\), with absolute multiplier
\(s_{l,j}(D_{l,j})_{i_{j+1},i_{j+1}}\).  Hence, with
\[
 \varepsilon_{l,i,j}
   =L_j^{-1}\log(D_{l,j})_{ii}-d_{l,i},
 \qquad \max_{l,i}|\varepsilon_{l,i,j}|\longrightarrow0
\]
by \eqref{eq:diagonal-growth}, one has
\begin{equation}
 \log\|Y(t_J)\|=
 \sum_{j=j_0}^{J-1}
 \left(\log s_{l,j}
       +j^7\bigl(d_{l,i_{j+1}}+\varepsilon_{l,i_{j+1},j}\bigr)\right).
 \label{eq:cycle-sum}
\end{equation}
Signs in \(\Pi_l\) do not affect this norm.

We verify the weighted average in \eqref{eq:cycle-sum}.  Group
the sum of \(j^7d_{l,i_{j+1}}\) into complete cycles of length
\(M_l\).  Within a cycle starting at \(j\),
\((j+r)^7-j^7=O(j^6)\) for \(0\leq r<M_l\), with a constant
depending on this fixed length.  Since the sum of
\(d_{l,i_{j+r+1}}-\overline d_l\) in a complete cycle is zero,
that cycle contributes only \(O(j^6)\) to the difference from
the mean-weighted sum.  Summing these errors and bounding the
at most \(M_l-1\) final terms by \(O(J^7)\) gives
\[
 \sum_{j=j_0}^{J-1}j^7d_{l,i_{j+1}}
 =\overline d_l\sum_{j=j_0}^{J-1}j^7+O(J^7).
\]
Also \(\sum_{j<J}|\log s_{l,j}|=O(\sum_{j<J}j^6)=O(J^7)\).
The \(\varepsilon\)-terms sum to \(o(J^8)\): for any positive
number, choose a fixed late index after which their absolute
values are below that number, and use
\(\sum_{j<J}j^7=O(J^8)\).  The finitely many earlier terms
vanish after division by \(J^8\).  Finally,
\[
 \sum_{j=j_0}^{J-1}j^7=\frac{J^8}{8}+O(J^7),\qquad
 t_J=\frac{J^8}{8}+O(J^7)
\]
by \eqref{eq:periods}.  Therefore
\begin{equation}
 \lim_{J\to\infty}\frac{\log\|Y(t_J)\|}{t_J}
 =\overline d_l<k.
 \label{eq:endpoint-rate}
\end{equation}

The estimate extends to every radius.  From \(Y'=P_lY\) and
\eqref{eq:riccati-barriers},
\[
 \|Y(t)\|\leq
 \exp\bigl(C_P(t-t_j)\bigr)\|Y(t_j)\|
 \qquad(t_j\leq t\leq t_{j+1}).
\]
The ratio \((t_{j+1}-t_j)/t_j\) tends to zero by
\eqref{eq:time-asymptotic}.  Thus \eqref{eq:endpoint-rate} gives
\[
 \limsup_{t\to\infty}\frac{\log\|Y(t)\|}{t}
 \leq\overline d_l.
\]
Choose \(\alpha_l\) with
\(\overline d_l<\alpha_l<k\).  For all sufficiently large
\(r=e^t\), \(\|Y(\log r)\|\leq r^{\alpha_l}\).
The fixed finite-dimensional space \(V_l\) has
\[
 \left|\sum_iY_i e_{l,i}(\omega)\right|\leq C_l\|Y\|
 \quad\text{for every }\omega\in S^m,
\]
for example by Cauchy--Schwarz and boundedness of the finitely
many smooth basis functions.  Since \(d_g(0,x)=r(x)\), this
gives \(|u(x)|\leq C_u(1+d_g(0,x))^k\) outside a compact set.
Smoothness bounds \(u\) on that compact set and supplies the
same inequality everywhere.

These functions are linearly independent: their restrictions
to the sphere \(t=t_{j_0}\) are the selected vectors of
mutually orthogonal round spaces \(V_l\).  It follows that
\[
 \dim_{\mathbb R}\mathcal H_k(\mathbb R^{m+1},g)
 \geq\sum_{l=2}^{L}M_l>\sum_{l=0}^{k}N_l.
\]
The Euclidean identification was proved at the end of
Section~\ref{sec:counting}, so this is the required strict comparison.

Finally, the distance and volume statements in Proposition~\ref{prop:metric} give
\[
 \operatorname{vol}_g B_g(0,R)
 =\operatorname{vol}_{g_0}(S^m)
       \int_0^R r^m a(\log r)^m\,dr.
\]
Because \(a(t)\to a_*\), changing variables \(r=Rz\) and
using dominated convergence shows that this quantity divided
by \(\omega_{m+1}R^{m+1}\) tends to \(a_*^m\); here
\(\operatorname{vol}_{g_0}(S^m)=(m+1)\omega_{m+1}\).
This completes the proof.
\end{proof}

\subsection{Cone limits, conformal curvature, and cone degrees}
\begin{proposition}[Cone geometry and the union of degrees]
\label{prop:cone-consequences}
Let \(g\) be the metric constructed in the proof of Theorem~\ref{thm:main},
with link parameters \(a_*,q_*\) and comparison integer \(k\).
Every cone with link \(G(a_*,q,J_0)\), \(1\leq q\leq q_*\), occurs
as a pointed tangent cone at infinity. In particular the tangent cone is
not unique. The metric \(g\) is not locally conformally flat. If
\(\mathcal D(M)\) denotes the union of nonnegative homogeneous harmonic
degrees over all its tangent cones, then \(\mathcal D(M)\) contains a
half-line and a neighborhood of \(k\).
\end{proposition}

\begin{proof}
 Centering a fixed
logarithmic window in the round rests or nonround dwells and letting its
center tend to infinity gives smooth annular limits with links
$G(a_*,1,J_0)$ and $G(a_*,q_*,J_0)$, respectively. They are nonisometric:
the first is Einstein, whereas the second has distinct horizontal and
vertical Ricci eigenvalues. Every intermediate ramp link also occurs as a
limit, because the ramp durations diverge and their derivatives tend to
zero. These annular limits extend to pointed cone limits including the
vertex: the inner radial ball of radius $\delta$ has diameter at most
$2\delta$ in every rescaling and in the limit. Comparing paths on the
complementary annuli, then letting $\delta\downarrow0$, gives convergence
of distances on each bounded ball. All these links have the same volume.

The additional local conformal flatness hypothesis in~\cite{CaiLai}
also fails.  Put \(h_*=G(a_*,q_*,J_0)\).
The nonround cone metric \(dr^2+r^2h_*\) is conformal, under
\(t=\log r\), to the product \(dt^2+h_*\).
For this product the Weyl tensor has mixed component
\[
 W_{0a0b}
 =-\frac{1}{m-1}
   \left(\Ric_{h_*}-\frac{\operatorname{Scal}_{h_*}}{m}h_*\right)_{ab}.
\]
Indeed the product has zero mixed curvature and radial Ricci tensor;
substitution in the trace decomposition of curvature gives this formula.
The component is nonzero because the link has distinct horizontal and
vertical Ricci eigenvalues.  Vanishing of the Weyl tensor is conformally
invariant and is preserved by smooth annular limits.  If the constructed
metric were locally conformally flat, its nonround cone limit would have
zero Weyl tensor, contradicting the formula.

For maximal Hopf charge in degree $l$, its
growth root $\alpha_l(q)$ is the nonnegative solution of
\[
 \alpha_l(q)(\alpha_l(q)+m-1)
 =a_*^{-2}q^{1/m}
   \bigl[l(l+m-1)+(q^{-1}-1)l^2\bigr].
\]
This root varies continuously with $q$. At $q=1$ its asymptotic slope is
$a_*^{-1}$; at $q=q_*$ it is
$a_*^{-1}q_*^{-(m-1)/(2m)}$, which is strictly smaller. More precisely, the endpoint roots are
\[
 \alpha_l(1)=a_*^{-1}l+O(1),\qquad
 \alpha_l(q_*)=a_*^{-1}q_*^{-(m-1)/(2m)}l+O(1).
\]
The difference of their slopes is positive, so for every sufficiently
large \(l\), \(\alpha_{l+1}(q_*)<\alpha_l(1)\) and
\(\alpha_l(q_*)<\alpha_l(1)\). The continuous image of the
\(q\)-interval contains the entire interval between these endpoints.
These intervals overlap and have unbounded upper endpoints, so their
union contains a half-line. In particular, equal cone-link
volumes do not provide a discrete union spectrum.

The chosen integer $k$ itself lies in this union. Indeed,
\eqref{eq:selection} forces some selected degree $l>k$. Its smallest dwell
root is below $k$, while its round root $\theta_l$ is greater than $l$
because $a_*<1$. Continuity supplies an intermediate cone root equal to
$k$; the strict endpoint inequalities in fact put a neighborhood of $k$
in $\mathcal D(M)$. Xu's nonunique-cone three-circles
theorem~\cite[Theorem~3.4]{Xu} requires the comparison exponent to lie
outside $\mathcal D(M)$, so it cannot be applied at $k$ or at any
sufficiently nearby exponent. The unique-cone spectral bounds in~\cite{Huang} also
have a hypothesis that fails. A global logarithmic growth average, as
in~\eqref{eq:endpoint-rate}, need not equal the local homogeneous degree
seen after normalization on each separate cone subsequence.
\end{proof}
